\hypertarget{appendix-ba8c1a9fe03f}{}
\section*{Default and value initialization}
\section{Default and value initialization}\label{default-and-value-initialization}

\textbf{Curriculum depth:} Introduced · \textbf{First seen:} C009 · \textbf{Latest development:} C009

\subsection{Definition and mechanism}\label{definition-and-mechanism}

No scalar initializer can leave automatic integers indeterminate. Aggregate empty braces initialize omitted members, whereas a user-provided empty default constructor may leave its scalar member uninitialized even with braces. First-declaration defaulting differs from out-of-class defaulting.

\subsection{Prerequisites}\label{prerequisites}

\begin{itemize}
\tightlist
\item
  \hyperref[appendix-7f8ec83d0886]{Constructors and invariant establishment} --- Initialization form and structural order determine how the constructor establishes members.
\end{itemize}

\subsection{Used by}\label{used-by}

\begin{itemize}
\tightlist
\item
  \hyperref[appendix-42e02a12c694]{Special member generation and traits}
\end{itemize}

\subsection{Recall and next depth}\label{recall-and-next-depth}

Classify int n, int n\{\}, defaulted A\{\}, and empty-body B\{\} without reading the bad cases. Pending: initializer-list overload selection and language-version uninitialized-value rules. Written coverage does not demonstrate personal mastery.

\subsection{Source and navigation}\label{source-and-navigation}

\href{Continuous\%20Notes/Computing\%20Foundations\%20-\%20From\%20Bits\%20to\%20C++\%20and\%20Systems\%20APIs\#Empty\%20syntax\%20does\%20not\%20promise\%20initialized\%20members}{Teaching mechanism and examples} · \hyperref[appendix-fd090d982a15]{Exact coverage} · \href{}{Dependency Map (02-Foundations-II, appendix-4ad47bdd269a)}
